\documentclass[11pt,a4paper]{article}

\usepackage[utf8]{inputenc}
\usepackage[T1]{fontenc}
\usepackage{amsmath,amssymb,amsthm}
\usepackage{mathtools}
\usepackage{bm}
\usepackage{graphicx}
\usepackage{booktabs}
\usepackage{hyperref}
\usepackage[margin=1in]{geometry}
\usepackage{algorithm}
\usepackage{algorithmic}
\usepackage{subcaption}
\usepackage{xcolor}
\usepackage{cleveref}

% Theorem environments
\newtheorem{definition}{Definition}
\newtheorem{proposition}{Proposition}
\newtheorem{theorem}{Theorem}
\newtheorem{remark}{Remark}

% Notation shortcuts
\newcommand{\field}{\bm{\phi}}
\newcommand{\deposit}{\bm{d}}
\newcommand{\evap}{\varepsilon}
\newcommand{\retain}{\alpha}
\newcommand{\R}{\mathbb{R}}
\newcommand{\E}{\mathbb{E}}
\newcommand{\x}{\bm{x}}
\newcommand{\q}{\bm{q}}
\newcommand{\K}{\bm{k}}
\newcommand{\vv}{\bm{v}}
\newcommand{\W}{W}

\title{Stigmergic Attention: Closing the Feedback Loop\\
Between Pheromone Fields and Neural Attention}

\author{Research Notes}
\date{\today}

\begin{document}
\maketitle

\begin{abstract}
We introduce \emph{stigmergic attention}, a modification of the transformer
attention mechanism that augments standard query-key matching with a
decaying pheromone field---an additive bias accumulated from past deposits
and subject to uniform erosion. Drawing on the theory of phase transitions
in lattice ant models, we identify a precise three-ingredient recipe for
edge-of-chaos dynamics: \emph{amplification} (positive feedback between
field and output), \emph{injection} (content-dependent deposits uncorrelated
with the field), and \emph{erosion} (uniform exponential decay). We show
that a na\"ive implementation lacking positive feedback reduces to an inert
decaying average with no phase transition, while closing the feedback
loop---computing deposits from attention \emph{output} rather than
input---recovers the full stigmergic dynamics. We present the architecture,
its connection to linear state-space models, and experimental results on
in-context learning tasks.
\end{abstract}

% ================================================================
\section{Introduction}
\label{sec:intro}

Stigmergy is the mechanism by which organisms coordinate through
modification of a shared environment rather than direct communication.
The canonical example is ant trail formation: ants deposit pheromone
as they walk, subsequent ants preferentially follow high-pheromone
paths, and pheromone evaporates uniformly over time. This creates a
positive feedback loop that, at the right balance between deposition
and evaporation, produces emergent structure at the edge of chaos
\cite{antpaper}.

Transformers \cite{vaswani2017attention} coordinate information flow
through attention---a content-based routing mechanism where each
position queries all others. We ask: can the stigmergic mechanism
augment attention with a complementary structural channel? Specifically,
can a decaying deposit field, accumulated causally across sequence
positions, provide useful inductive bias for in-context learning?

Our investigation proceeds in three stages:
\begin{enumerate}
    \item We define stigmergic attention as standard attention plus an
          additive field bias from a linear recurrence (\Cref{sec:architecture}).
    \item We identify, via the ant trail phase transition literature, that
          positive feedback is the critical missing ingredient, and show
          how to close the loop (\Cref{sec:theory}).
    \item We present experiments comparing feedback and non-feedback
          variants across evaporation rates (\Cref{sec:experiments}).
\end{enumerate}


% ================================================================
\section{Architecture}
\label{sec:architecture}

% ----------------------------------------------------------------
\subsection{Standard Attention}

Given an input sequence $X \in \R^{T \times d}$, standard causal
multi-head attention computes, for each head $h \in \{1, \ldots, H\}$:

\begin{definition}[Scaled dot-product attention]
\label{def:attention}
Let $\q_i, \K_j \in \R^{d_h}$ be the query and key vectors at positions
$i$ and $j$ respectively, with $d_h = d / H$ the head dimension. The
attention weight from position $i$ to position $j$ is:
\begin{equation}
    a_{ij}^{(h)} = \frac{\exp\!\bigl(\q_i^{(h)\top} \K_j^{(h)} / \sqrt{d_h}\bigr)}
                        {\sum_{k \leq i} \exp\!\bigl(\q_i^{(h)\top} \K_k^{(h)} / \sqrt{d_h}\bigr)}
    \quad \text{for } j \leq i
\end{equation}
where the restriction $j \leq i$ enforces causal masking.
The output at position $i$ is:
\begin{equation}
    \bm{o}_i = \text{Concat}_{h=1}^{H}\!\Bigl(\sum_{j \leq i} a_{ij}^{(h)} \vv_j^{(h)}\Bigr) \cdot W_O
\end{equation}
with queries $\q^{(h)} = X W_Q^{(h)}$, keys $\K^{(h)} = X W_K^{(h)}$,
values $\vv^{(h)} = X W_V^{(h)}$, and output projection $W_O \in \R^{d \times d}$.
\end{definition}

% ----------------------------------------------------------------
\subsection{The Pheromone Field}

We augment attention with a scalar field per head that accumulates
deposits and decays over positions.

\begin{definition}[Pheromone field]
\label{def:field}
For head $h$, the \emph{pheromone field} at position $i$ is defined
by the linear recurrence:
\begin{equation}
    \phi_i^{(h)} = \retain \cdot \phi_{i-1}^{(h)} + d_i^{(h)},
    \quad \phi_0^{(h)} = d_0^{(h)}
    \label{eq:field_recurrence}
\end{equation}
where:
\begin{itemize}
    \item $\retain = 1 - \evap \in (0, 1)$ is the \emph{retention rate},
          with $\evap$ the \emph{evaporation rate}
    \item $d_i^{(h)} \in \R$ is the \emph{deposit} at position $i$ for head $h$
    \item The field has initial condition $\phi_{-1}^{(h)} = 0$
\end{itemize}
\end{definition}

\begin{proposition}[Closed-form solution]
\label{prop:closedform}
The recurrence in \Cref{eq:field_recurrence} has the closed-form solution:
\begin{equation}
    \phi_i^{(h)} = \sum_{j=0}^{i} \retain^{i-j} \cdot d_j^{(h)}
    \label{eq:field_closedform}
\end{equation}
which can be computed as a matrix-vector product $\field^{(h)} = D \deposit^{(h)}$
where $D \in \R^{T \times T}$ is the lower-triangular decay matrix with
entries $D_{ij} = \retain^{i-j}$ for $j \leq i$ and $0$ otherwise.
\end{proposition}

\begin{proof}
By induction. Base case: $\phi_0 = d_0 = \retain^0 d_0$. Inductive step:
$\phi_i = \retain \sum_{j=0}^{i-1} \retain^{i-1-j} d_j + d_i
= \sum_{j=0}^{i-1} \retain^{i-j} d_j + \retain^0 d_i
= \sum_{j=0}^{i} \retain^{i-j} d_j$.
\end{proof}

\begin{remark}[Connection to state-space models]
\Cref{eq:field_recurrence} is a first-order linear recurrence---the
simplest possible state-space model (SSM). It is a scalar, per-head
instance of the general form $s_i = A s_{i-1} + B u_i$ with
$A = \retain \in \R$ and $B = 1$. Unlike S4 or Mamba, which use
structured matrices $A \in \R^{n \times n}$ for multi-dimensional
state, the pheromone field uses a single scalar state per head,
making it both interpretable and cheap.
\end{remark}

% ----------------------------------------------------------------
\subsection{Stigmergic Attention}

The field enters the attention computation as an additive bias on the
pre-softmax logits.

\begin{definition}[Stigmergic attention]
\label{def:stigmergic}
The stigmergic attention logit from position $i$ to position $j$ is:
\begin{equation}
    \ell_{ij}^{(h)} = \frac{\q_i^{(h)\top} \K_j^{(h)}}{\sqrt{d_h}} + \phi_j^{(h)}
    \label{eq:stigmergic_logit}
\end{equation}
The attention weights are obtained by applying causal masking and softmax
over the logit vector $\bm{\ell}_i^{(h)} = [\ell_{i1}^{(h)}, \ldots, \ell_{iT}^{(h)}]$
as in \Cref{def:attention}.
\end{definition}

The field bias $\phi_j^{(h)}$ is \emph{position-indexed}: it biases
attention \emph{toward} position $j$ across all queries $i \geq j$.
Positions with high accumulated deposits attract more attention,
analogous to how pheromone-rich trail segments attract more ants.

The key design question is: where do deposits come from?


% ================================================================
\section{Theory: The Three-Ingredient Recipe}
\label{sec:theory}

% ----------------------------------------------------------------
\subsection{Lessons from Ant Trail Phase Transitions}

The lattice ant model \cite{antpaper} studies pheromone trail formation
on a discrete grid. The continuum limit of the pheromone dynamics is:
\begin{equation}
    \frac{\partial \phi}{\partial t}
    = \underbrace{\alpha_{\text{inj}}}_{\text{injection}}
    + \underbrace{\beta \cdot g(\phi)}_{\text{amplification}}
    - \underbrace{\lambda \cdot \phi}_{\text{erosion}}
    + \underbrace{D \nabla^2 \phi}_{\text{diffusion}}
    \label{eq:ant_pde}
\end{equation}
The paper's central finding is that edge-of-chaos dynamics---the boundary
between frozen (ordered) and disordered phases---requires exactly three
ingredients, all \emph{structurally uncorrelated} with the field $\phi$:

\begin{enumerate}
    \item \textbf{Amplification}: Positive feedback where the field
          influences behavior, which in turn reinforces the field.
          In ants: high-pheromone trails attract more ants, who deposit
          more pheromone.
    \item \textbf{Injection}: An unconditional source of deposits that
          is uncorrelated with the current field state. In ants: all ants
          deposit pheromone regardless of the existing trail strength.
    \item \textbf{Erosion}: Uniform removal that is uncorrelated with the
          field state. In ants: pheromone evaporates at a constant rate
          everywhere.
\end{enumerate}

The paper tested seven alternative mechanisms that replace uniform erosion
with state-dependent processes (habituation, peaked response functions,
Hill functions, coupled deposition). \emph{Every alternative failed} to
produce edge-of-chaos dynamics, because state-dependent disorder is
correlated with the field---it preferentially targets the structures that
contribute most to the order parameter.

% ----------------------------------------------------------------
\subsection{Mapping to Stigmergic Attention}

We now map the three ingredients to the neural architecture:

\begin{table}[h]
\centering
\begin{tabular}{@{}llll@{}}
\toprule
\textbf{Ingredient} & \textbf{Ant system} & \textbf{Neural analogue} & \textbf{Property} \\
\midrule
Amplification & Ants follow pheromone, & Field biases attention, & Positive feedback \\
              & deposit more on trails & output generates deposits & ($\phi$-correlated) \\
\addlinespace
Injection     & All ants deposit & Token content produces & Uncorrelated \\
              & unconditionally & deposits via $W_d$ & with $\phi$ \\
\addlinespace
Erosion       & Uniform evaporation & Exponential decay & Uncorrelated \\
              & at rate $\lambda$ & at rate $\evap$ per position & with $\phi$ \\
\bottomrule
\end{tabular}
\caption{Three-ingredient mapping from ant systems to stigmergic attention.}
\label{tab:ingredients}
\end{table}

The critical question is whether the architecture implements
\emph{amplification}---positive feedback between the field and the
deposit process.

% ----------------------------------------------------------------
\subsection{The Non-Feedback Architecture (Negative Result)}
\label{sec:nonfeedback}

In the na\"ive implementation, deposits are computed from the
\emph{input} to the attention layer:
\begin{equation}
    d_i^{(h)} = W_d^{(h)} \x_i
    \label{eq:deposit_input}
\end{equation}
where $W_d^{(h)} \in \R^{1 \times d}$ is a learned projection and
$\x_i \in \R^d$ is the input embedding at position $i$.

\begin{proposition}[No positive feedback without output deposits]
\label{prop:nofeedback}
When deposits are computed from inputs via \Cref{eq:deposit_input},
the deposit $d_i^{(h)}$ is independent of the field
$\phi_0^{(h)}, \ldots, \phi_{i-1}^{(h)}$. The field is therefore
a deterministic function of the input sequence alone:
\begin{equation}
    \phi_i^{(h)} = \sum_{j=0}^{i} \retain^{i-j} \cdot W_d^{(h)} \x_j
\end{equation}
This is a causally-weighted moving average of input projections---an
exponentially-weighted positional encoding with no feedback mechanism.
\end{proposition}

Without amplification, the field cannot create or maintain
structure beyond what the input sequence provides. No phase
transition is possible because there is no mechanism for
self-reinforcement.

% ----------------------------------------------------------------
\subsection{Closing the Feedback Loop}
\label{sec:feedback}

The fix is to compute deposits from the attention \emph{output}:

\begin{definition}[Feedback stigmergic attention]
\label{def:feedback}
At each position $i$, the computation proceeds causally:
\begin{enumerate}
    \item Compute field bias from accumulated past deposits:
    \begin{equation}
        \phi_i^{(h)} = \retain \cdot \phi_{i-1}^{(h)} + d_{i-1}^{(h)}
    \end{equation}
    \item Compute attention weights with field bias:
    \begin{equation}
        a_{ij}^{(h)} = \text{softmax}_j\!\Bigl(
            \frac{\q_i^{(h)\top} \K_j^{(h)}}{\sqrt{d_h}} + \phi_j^{(h)}
        \Bigr) \quad \text{for } j \leq i
    \end{equation}
    \item Compute attention output at position $i$:
    \begin{equation}
        \bm{o}_i^{(h)} = \sum_{j \leq i} a_{ij}^{(h)} \vv_j^{(h)}
    \end{equation}
    \item Compute deposit from output:
    \begin{equation}
        d_i^{(h)} = W_d^{(h)} \bm{o}_i
        \label{eq:deposit_output}
    \end{equation}
    where $\bm{o}_i = \text{Concat}_h(\bm{o}_i^{(h)}) \cdot W_O$.
\end{enumerate}
\end{definition}

This creates the closed feedback loop:
\begin{equation}
    \phi \xrightarrow{\text{bias}} \text{attention}
    \xrightarrow{\text{aggregate}} \text{output}
    \xrightarrow{W_d} \text{deposit}
    \xrightarrow{\text{accumulate}} \phi
    \label{eq:feedback_loop}
\end{equation}

\begin{remark}[Causality]
\label{rem:causality}
Position $i$'s deposit $d_i^{(h)}$ depends on the field at positions
$0, \ldots, i-1$ (through the attention bias) and on the values at
positions $0, \ldots, i$ (through the attention output). The deposit
at position $i$ only affects the field at positions $i+1, i+2, \ldots$.
This maintains strict causality: no position can influence its own
field bias, analogous to how an ant's deposit only affects future ants.
\end{remark}

\begin{remark}[Computational cost]
The feedback architecture requires sequential processing over positions
(an $O(T)$ loop) because each position's output depends on the field
state, which depends on all previous outputs. In parallel mode
(training), this can be amortized: the field recurrence is a single
scan operation, and the attention computation at each step only
requires the pre-computed $Q, K, V$ projections. The overhead per
position is a single linear projection $W_d$ (size $d \to H$) and
a scalar multiply-add per head for the field update.
\end{remark}


% ================================================================
\section{Analysis}
\label{sec:analysis}

% ----------------------------------------------------------------
\subsection{The Field as Modulated State-Space Model}

In the non-feedback variant, the pheromone field is a standard linear
SSM with fixed dynamics. In the feedback variant, the effective dynamics
become input-dependent:

\begin{proposition}[Feedback field as nonlinear SSM]
\label{prop:nonlinear}
The feedback field can be written as:
\begin{equation}
    \phi_i = \retain \cdot \phi_{i-1} + f_\theta(\x_{0:i}, \phi_{0:i-1})
\end{equation}
where $f_\theta$ is the composition of attention (with field-dependent
weights) and the deposit projection. This is a nonlinear SSM: the
transition function depends on both the input sequence and the state
history through the attention mechanism.
\end{proposition}

This is qualitatively different from linear SSMs (S4, Mamba): the
state dynamics are \emph{modulated} by the attention mechanism, which
itself is modulated by the state. The field and attention co-evolve.

% ----------------------------------------------------------------
\subsection{Information-Theoretic Interpretation}

The three ingredients admit an information-theoretic reading:

\begin{itemize}
    \item \textbf{Erosion} ($\retain < 1$): Implements an information
          horizon. The contribution of position $j$ to the field at
          position $i$ decays as $\retain^{i-j}$, creating an
          exponential forgetting curve. The effective memory span is
          $\tau_{\text{eff}} = -1/\!\ln(\retain) \approx 1/\evap$
          for small $\evap$.

    \item \textbf{Injection}: Ensures the field receives fresh signal
          at every position, preventing collapse to a fixed point.
          The deposit $d_i$ injects new information proportional to the
          local content, maintaining exploration of the input space.

    \item \textbf{Amplification}: Creates selective reinforcement.
          Positions that receive high attention produce outputs that
          generate deposits that increase future attention to
          positions with similar patterns. This is a form of
          \emph{iterated distillation}: the field accumulates a
          compressed summary of what the attention mechanism has found
          useful.
\end{itemize}

The evaporation rate $\evap$ controls the balance:
\begin{itemize}
    \item $\evap \to 0$ (frozen): infinite memory, field dominated by
          early deposits, attention locked to initial patterns
    \item $\evap \to 1$ (disordered): instant forgetting, field
          reflects only the most recent deposit, no accumulation
    \item $\evap$ at criticality: memory span matches the task's
          relevant context window
\end{itemize}


% ----------------------------------------------------------------
\subsection{Relationship to Existing Architectures}

\begin{table}[h]
\centering
\begin{tabular}{@{}lcccc@{}}
\toprule
\textbf{Architecture} & \textbf{State dim.} & \textbf{Feedback} &
\textbf{Attention} & \textbf{Decay} \\
\midrule
Standard Transformer & --- & --- & content-based & --- \\
ALiBi               & --- & --- & content + static bias & --- \\
Feedback Transformer & $d$ & nonlinear & content-based & --- \\
RetNet              & $d$ & linear & via $\gamma$-mask & multiplicative \\
S4 / Mamba          & $n$ & linear & selective (Mamba) & structured $A$ \\
Linear Attention    & $d_k \times d_v$ & linear & via kernel & implicit \\
\addlinespace
Stigmergic (input)  & $H$ & none & content + field bias & additive \\
Stigmergic (feedback) & $H$ & nonlinear & content + field bias & additive \\
\bottomrule
\end{tabular}
\caption{Comparison of recurrent/attention architectures. $H$ = number
of heads, $d$ = model dimension, $n$ = SSM state dimension.}
\label{tab:architectures}
\end{table}

Stigmergic attention decomposes into three familiar components---a
scalar recurrence, an additive logit bias, and output feedback---each
of which appears independently in prior work.  The contribution is
their specific combination, motivated by the three-ingredient theory
rather than by architectural search.

\paragraph{Additive logit biases.}
ALiBi \cite{press2022alibi} adds a fixed, non-learned linear penalty
to attention logits based on key--query distance, improving length
generalization without positional embeddings.  T5 relative position
biases \cite{raffel2020t5} learn a per-distance scalar bias.  Both
are \emph{static}: the bias depends only on position, not on content
or state.  The pheromone field replaces these fixed biases with a
content-dependent recurrence---it is ``ALiBi where the slope is a
running state driven by what attention has found useful.''

\paragraph{Output feedback.}
The Feedback Transformer \cite{fan2021feedback} feeds the full
$d$-dimensional output of the top layer back as additional context
for the next position's attention computation.  This implements
the same principle of output-dependent state, but through a
full-dimensional vector used as an extra key--value pair rather than
a scalar logit bias.  The stigmergic field is a more constrained
version: a single scalar per head, which makes the recurrence
interpretable and cheap but limits expressivity.

\paragraph{SSM--attention hybrids.}
Jamba \cite{lieber2024jamba}, Griffin \cite{de2024griffin}, and
RWKV \cite{peng2023rwkv} interleave SSM or linear-recurrence layers
with attention layers.  In all cases the two mechanisms are
\emph{decoupled}: the SSM layer processes its input independently,
then passes its output to attention (or vice versa).  There is no
feedback loop between the recurrent state and the attention weights.
Mamba-2 \cite{dao2024mamba2} shows a formal duality between SSMs and
structured attention, but replaces softmax attention with a linear
variant rather than augmenting it.

\paragraph{What is novel.}
The specific wiring---a scalar recurrence whose output additively
biases softmax attention logits, with the recurrence input derived
from attention output---does not appear in prior work to our
knowledge.  However, the individual components are not new, and a
reviewer could fairly characterise the architecture as ``ALiBi with a
content-dependent recurrence'' or ``a 1-d Mamba biasing attention
logits.''  We do not claim the architecture itself as the primary
contribution.  Rather, the value lies in the \emph{theoretical
framework}: the three-ingredient recipe from ant trail phase
transitions provides principled answers to design questions that
would otherwise be arbitrary---why scalar (not vector)? why additive
(not multiplicative)? why from output (not input)? why uniform decay
(not state-dependent)?  Without this framework, the non-feedback
variant (\Cref{sec:nonfeedback}) looks equally reasonable, and there
is no theoretical reason to predict its failure.


% ================================================================
\section{Experiments}
\label{sec:experiments}

% ----------------------------------------------------------------
\subsection{Setup}

\paragraph{Model.}
A micro-GPT with 2 layers, 4 heads, $d_{\text{model}} = 48$,
ReLU activation, RMSNorm, and $\sim$58K parameters. This deliberately
small scale allows exhaustive sweeps with multiple seeds.

\paragraph{Task.}
In-context modular arithmetic: each sequence presents 8 input-output
pairs $(x_i, y_i)$ where $y_i = (ax_i + b) \bmod V$ for a random
function $(a, b)$, followed by a query $x_q$. The model must predict
$y_q$. Vocabulary $V = 16$, sequence length $T = 18$.
Training and test function sets are disjoint (200 train / 40 test
functions), ensuring evaluation measures out-of-distribution
generalization.

\paragraph{Training.}
AdamW optimizer, learning rate $3 \times 10^{-3}$ with linear decay,
weight decay 0.10, dropout 0.05, batch size 64, 5000 steps.
Validation accuracy averaged over last 5 checkpoints (every 50 steps).

% ----------------------------------------------------------------
\subsection{Non-Feedback Field (Negative Result)}

We first sweep evaporation rate $\evap \in [0, 0.40]$ with the
non-feedback architecture (deposits from input, \Cref{eq:deposit_input}).

\begin{table}[h]
\centering
\begin{tabular}{@{}lccc@{}}
\toprule
$\evap$ & Val. accuracy & $\pm$ std & Seeds \\
\midrule
0.00 & 0.656 & 0.015 & 3 \\
0.05 & 0.652 & 0.012 & 3 \\
0.10 & 0.653 & 0.018 & 3 \\
0.15 & 0.648 & 0.011 & 3 \\
0.20 & 0.674 & 0.009 & 3 \\
0.25 & 0.663 & 0.011 & 3 \\
0.30 & 0.651 & 0.012 & 3 \\
0.40 & 0.722 & 0.089 & 3 \\
\addlinespace
No field & 0.722 & --- & 3 \\
\bottomrule
\end{tabular}
\caption{Non-feedback (decay) field: no $\evap$ value matches the
no-field baseline. The field is inert.}
\label{tab:decay_results}
\end{table}

As predicted by \Cref{prop:nofeedback}, the non-feedback field provides
no benefit. All $\evap$ values produce validation accuracy at or below
the no-field baseline. There is no phase transition---the field is a
harmless but useless exponentially-weighted positional encoding.

% ----------------------------------------------------------------
\subsection{Feedback Field}

We repeat the sweep with the feedback architecture (deposits from
output, \Cref{eq:deposit_output}).

\begin{table}[h]
\centering
\begin{tabular}{@{}lccc@{}}
\toprule
$\evap$ & Val. accuracy & $\pm$ std & Seeds \\
\midrule
0.00 & 0.654 & 0.016 & 3 \\
0.05 & \textbf{0.715} & 0.080 & 3 \\
0.10 & 0.704 & 0.069 & 3 \\
0.15 & 0.654 & 0.015 & 3 \\
0.20 & 0.662 & 0.009 & 3 \\
0.25 & 0.713 & 0.086 & 3 \\
0.30 & 0.697 & 0.074 & 3 \\
0.40 & 0.647 & 0.018 & 3 \\
\addlinespace
No field & 0.722 & 0.069 & 3 \\
\bottomrule
\end{tabular}
\caption{Feedback field results: deposits computed from attention output.
The feedback mode shows higher variance and seed-dependent bimodality
(one seed consistently reaches $\sim$0.80--0.83 while others stay at
$\sim$0.65), suggestive of a phase transition that the model can reach
but does not reliably find.}
\label{tab:feedback_results}
\end{table}

The feedback results are nuanced. The mean accuracy does not exceed
the no-field baseline, but the \emph{variance structure} is qualitatively
different from the decay mode. In the decay mode, all seeds cluster
tightly around 0.65. In the feedback mode, seed 137 consistently reaches
0.80--0.83 across multiple $\evap$ values while other seeds remain at
$\sim$0.65. This bimodality is characteristic of systems near a phase
transition: the feedback loop creates a basin of attraction that some
initializations find and others miss.

\begin{remark}[Seed-dependent bimodality]
The high-variance points ($\evap = 0.05, 0.10, 0.25, 0.30$) all show
the same pattern: one seed near 0.80, two seeds near 0.65. This is
not random noise---it is consistent across $\evap$ values for the same
seed. The feedback loop appears to create a useful attractor that
requires specific initialization conditions to reach, suggesting that
the optimization landscape has multiple basins separated by barriers.
\end{remark}

% ----------------------------------------------------------------
\subsection{Length Generalization on State-Tracking Tasks}
\label{sec:lengthgen}

The modular arithmetic task is fundamentally about random-access
pattern matching---the model looks at examples, spots the function, and
applies it. Standard attention already excels at this. A more principled
evaluation tests the field on tasks requiring \emph{sequential state
accumulation}, evaluated on \emph{out-of-distribution sequence lengths}.

We train on sequences of length $T = 34$ ($n = 16$ examples) and test
on $T \in \{34, 66, 98, 130\}$, with three tasks:

\begin{itemize}
    \item \textbf{Parity}: given random tokens $x_0, \ldots, x_n \in
          [0, V\!-\!1]$, predict the running parity
          $(x_0 + \cdots + x_i) \bmod 2$ at each position. This is a
          classic transformer limitation: parity requires aggregating
          over \emph{all} previous inputs and is not reducible to
          attending to a single position.
    \item \textbf{Flip-flop}: a register machine where token $0$ is
          a no-op and any other token $k$ sets the state to $k$. The
          model must predict the current state---i.e., the last non-zero
          write, which may be arbitrarily far back. This tests
          long-range memory retention.
    \item \textbf{Cumulative sum}: predict $(x_0 + \cdots + x_i)
          \bmod V$ at each position. This requires precise modular
          arithmetic, not just binary state.
\end{itemize}

All tasks are trivially solvable with a running state variable (a
single register updated at each step) and all models achieve 1.000
accuracy at the training length $T = 34$. The question is whether
this capability \emph{generalizes} to unseen lengths.

\begin{table}[h]
\centering
\begin{tabular}{@{}lcccc@{}}
\toprule
& $T = 34$ & $T = 66$ & $T = 98$ & $T = 130$ \\
& (train) & (OOD) & (OOD) & (OOD) \\
\midrule
\multicolumn{5}{l}{\textbf{Parity}} \\
\quad No field       & 1.000 & 0.542{\tiny $\pm$.021} & 0.512{\tiny $\pm$.010} & 0.514{\tiny $\pm$.013} \\
\quad Feedback $\evap\!=\!0.05$ & 1.000 & \textbf{0.615}{\tiny $\pm$.065} & 0.526{\tiny $\pm$.029} & \textbf{0.554}{\tiny $\pm$.058} \\
\quad Feedback $\evap\!=\!0.10$ & 1.000 & 0.540{\tiny $\pm$.045} & \textbf{0.536}{\tiny $\pm$.059} & 0.525{\tiny $\pm$.022} \\
\quad Feedback $\evap\!=\!0.20$ & 1.000 & 0.540{\tiny $\pm$.045} & 0.535{\tiny $\pm$.049} & 0.521{\tiny $\pm$.013} \\
\addlinespace
\multicolumn{5}{l}{\textbf{Flip-flop}} \\
\quad No field       & 1.000 & 0.524{\tiny $\pm$.007} & 0.516{\tiny $\pm$.002} & 0.521{\tiny $\pm$.003} \\
\quad Feedback $\evap\!=\!0.05$ & 1.000 & \textbf{0.557}{\tiny $\pm$.011} & 0.518{\tiny $\pm$.004} & 0.524{\tiny $\pm$.009} \\
\addlinespace
\multicolumn{5}{l}{\textbf{Cumulative sum}} \\
\quad No field       & 1.000 & 0.098{\tiny $\pm$.013} & 0.072{\tiny $\pm$.012} & 0.078{\tiny $\pm$.014} \\
\quad Feedback $\evap\!=\!0.05$ & 1.000 & 0.083{\tiny $\pm$.012} & 0.076{\tiny $\pm$.027} & 0.054{\tiny $\pm$.008} \\
\bottomrule
\end{tabular}
\caption{Length generalization: train on $T=34$, test on longer
sequences. Random baseline is 0.5 for parity, $1/V \approx 0.063$ for
cumsum, $\sim 1/V$ for flip-flop. Bold marks the best result per
row where it exceeds no-field by $>$1$\sigma$. Three seeds per condition.}
\label{tab:lengthgen}
\end{table}

The results reveal a clear task-dependent pattern:

\paragraph{Parity (positive result).}
The feedback field at $\evap = 0.05$ improves OOD accuracy from
$0.542 \to 0.615$ at $T = 66$ and from $0.514 \to 0.554$ at
$T = 130$. The baseline drops to random (0.50) immediately outside
the training distribution, while the feedback field retains partial
generalization. This is the first clean positive result: the scalar
recurrence carries parity state---literally a running XOR---beyond
what position-tied attention patterns can represent.

\paragraph{Flip-flop (modest signal).}
Feedback $\evap = 0.05$ shows a consistent improvement at $T = 66$
($0.557$ vs $0.524$) but the effect fades at longer lengths. The
flip-flop task requires \emph{selective overwriting}---the field must
learn to replace its state when a non-zero token appears and maintain
it through no-ops. A single scalar recurrence with fixed decay rate
is a poor fit: the decay is unconditional, so the state fades even
when it should persist.

\paragraph{Cumulative sum (negative result).}
Both modes collapse to near-random at $T = 66$. Cumulative sum mod
$V = 16$ requires precise 4-bit modular arithmetic---a scalar field
with continuous values cannot faithfully encode a discrete counter.
The task demands \emph{multi-bit state} that exceeds the capacity of
$H$ scalar channels.

\paragraph{Why parity works and cumsum does not.}
The key distinction is \emph{state complexity}. Parity has 1 bit of
state (even/odd). A single scalar field head can encode this as the
sign of the field value, and the sign is preserved under the
exponential decay (positive stays positive, negative stays negative).
Cumulative sum mod 16 has $\lceil \log_2 16 \rceil = 4$ bits of
state. To encode this, the field would need to maintain 4 independent
channels that evolve without interference---exactly what the $H = 4$
heads could in principle provide, but the learned positional embeddings
(which do not generalize to unseen positions) destroy the field's
contribution at OOD lengths.

This suggests a clear scaling law for stigmergic attention: the
feedback field can track $O(H)$ bits of state, and should outperform
standard attention on tasks whose state complexity is $\leq H$ bits,
evaluated at lengths beyond the training distribution.


% ================================================================
\section{Discussion}
\label{sec:discussion}

\subsection{What the Bug Revealed}

Our investigation began with a serendipitous bug: an inverted decay
exponent that caused the decay matrix to \emph{amplify} old deposits
($D_{ij} = \retain^{j-i}$ with $j < i$, giving $D_{ij} > 1$) rather
than decay them. This buggy ``amplify'' mode produced a striking phase
transition---but one driven by exponential growth rather than
stigmergic dynamics. Fixing the bug eliminated the phase transition
entirely, revealing that the correct-decay field was inert.

The ant paper provided the explanation: the architecture had injection
(content-dependent deposits) and erosion (uniform decay) but lacked
\emph{amplification} (positive feedback). The buggy amplification
accidentally provided a crude substitute---not through the stigmergic
mechanism, but through raw exponential growth of the field values.

\subsection{The Feedback Hypothesis}

The feedback architecture (\Cref{def:feedback}) is the minimal change
that closes the loop. It adds no parameters (same $W_d$ projection)
and changes only the \emph{source} of deposits: from the attention
input $\x_i$ to the attention output $\bm{o}_i$.

This creates qualitatively different dynamics:
\begin{itemize}
    \item \textbf{Self-reinforcement}: positions that receive high
          attention produce outputs that generate deposits, which
          increase future attention to those positions.
    \item \textbf{Trail formation}: repeatedly useful patterns create
          persistent field structures that guide future attention.
    \item \textbf{Competition}: uniform erosion ensures that only
          actively reinforced structures persist---abandoned trails
          fade.
\end{itemize}

\subsection{Limitations}

\begin{itemize}
    \item \textbf{Scale}: All experiments use a $\sim$58K parameter
          model on a single synthetic task. The feedback mechanism's
          value at transformer scale is unknown.
    \item \textbf{Sequential bottleneck}: The feedback loop requires
          position-by-position processing during training, losing the
          parallelism advantage of standard attention. For long
          sequences, this may be prohibitive without approximation.
    \item \textbf{Task specificity}: In-context modular arithmetic is
          a controlled but narrow evaluation. Natural language tasks
          may not benefit from the same mechanism.
\end{itemize}


% ================================================================
\section{Research Directions}
\label{sec:directions}

% ----------------------------------------------------------------
\subsection{Finding the Right Task}
\label{sec:right_task}

The modular arithmetic task may be the wrong testbed. It is
fundamentally about \emph{algebraic structure}---recognizing $(a, b)$
from examples and applying a formula. Attention is already good at
this: match the query $x$ against example $x$'s, retrieve the
corresponding $y$, infer the pattern.

A stigmergic field should shine on tasks where \textbf{sequential
accumulation matters}---where the answer depends not on retrieving a
single example but on building up a representation \emph{across} the
sequence. Three candidate task families:

\begin{itemize}
    \item \textbf{Counting and frequency estimation}: ``Which token
          appeared most often?'' requires accumulating evidence across
          all positions, not just attending to one. The field naturally
          accumulates; attention does not.

    \item \textbf{Order-dependent state tracking}: ``Follow these
          instructions in sequence''---where each step modifies a
          state. A maze walk, a state machine, a sequence of edits.
          The field's recurrence is literally a running state.

    \item \textbf{Multi-hop reasoning with fading evidence}: ``A is
          north of B, B is north of C, where is A relative to C?''
          spread across many distractor tokens. The field can maintain
          a decaying trace of B's location while processing
          distractors, bridging the gap without needing to attend all
          the way back.
\end{itemize}

The common thread: tasks where a \textbf{running summary} is more
useful than \textbf{random-access retrieval}. Standard attention
excels at random access (query any position). The field adds a
complementary channel---a lossy, compressed, sequential summary. If
the task does not need that, the field is dead weight.

The modular arithmetic task is almost pure random access: look at the
examples, spot the pattern, apply it. No wonder the field does not
help.

A concrete proposal: a \textbf{state-tracking task} such as ``given a
sequence of $+1$ and $-1$ tokens, is the running sum positive or
negative at the query position?'' This is trivially easy with a
running state and provably hard for bounded-depth transformers. If
the feedback field does not help \emph{there}, it is dead. If it
does, we have found the right problem class.

% ----------------------------------------------------------------
\subsection{Alternative Architectural Framings}
\label{sec:alt_arch}

The feedback loop implements a general pattern:
\begin{equation*}
    \text{state} \to \text{modulate processing} \to \text{observe result}
    \to \text{update state} \to \text{repeat}
\end{equation*}
This pattern appears in several architectural families beyond
attention, suggesting alternative carriers for the stigmergic
mechanism.

\paragraph{Energy-based models.}
The field can be interpreted as an energy landscape that attention
navigates. Deposits lower energy at useful positions, erosion raises
it everywhere, and attention follows the energy gradient. This maps
to Hopfield networks---modern Hopfield layers \cite{ramsauer2021hopfield}
already connect energy models to attention. The feedback field could
be an \emph{asymmetric} Hopfield energy that evolves causally.

\paragraph{Predictive coding.}
Under the predictive coding hypothesis, each layer maintains a
prediction of its input, computes prediction errors, and updates.
The field serves as the prediction (running estimate of what is
important), the attention output is the observation, and the deposit
is the prediction error. This framing naturally explains why deposits
should come from the \emph{output} not the \emph{input}---the
deposit is the \emph{surprise}, not the stimulus.

\paragraph{Recurrent attention.}
Universal Transformers and PonderNet iterate the same attention layer
multiple times. The feedback field is a lightweight version: instead
of rerunning full attention, it carries forward a scalar summary per
head. Much cheaper, but the same principle of iterative refinement.

\paragraph{Diffusion models.}
In diffusion models, a noisy signal is iteratively refined by
predicting and removing noise. The pheromone field could be seen as a
``structure estimate'' that gets refined position by position. However,
diffusion operates in continuous space with a learned score function,
while our field is a scalar per head. The analogy is suggestive but
the mapping is loose.

We note that the hypothesis is about the \emph{feedback mechanism},
not the specific carrier. If positive feedback between a decaying
state and a processing layer produces useful dynamics, the principle
should transfer across architectures. We therefore resist switching
carriers until the mechanism is validated on a suitable task.

% ----------------------------------------------------------------
\subsection{Theoretical Limitations of Transformers}
\label{sec:limitations}

The strongest motivation for stigmergic attention comes from known
theoretical limitations of transformers that a feedback field directly
addresses.

\paragraph{State tracking.}
Transformers are provably bad at simulating automata or tracking
state over long sequences \cite{merrill2023transformers}. They cannot
count, cannot maintain a running tally, cannot simulate a state
machine beyond a few steps. The field IS a running state---a scalar
recurrence per head. It does not give a full state machine, but it
provides $H$ scalar channels of persistent state that attention can
read and write. That is exactly what is missing.

\paragraph{Length generalization.}
Transformers trained on length $N$ fail at length $N + k$. This is
because attention patterns are tied to absolute or relative positions.
The field's decay creates a \emph{scale-invariant} memory---the
exponential decay $\retain^{i-j}$ depends only on distance, not
absolute position. A feedback field trained on short sequences might
generalize to longer ones because the dynamics are
position-independent.

\paragraph{Compositional generalization.}
Multi-step reasoning requires chaining intermediate results. Attention
can only do this by allocating positions as ``scratchpad''
(chain-of-thought). The field provides an implicit scratchpad---each
position's output modifies the field, carrying forward a compressed
result without consuming a token position.

\medskip

The clearest theoretical pitch: \textbf{transformers are memoryless
between positions}---each position's computation is independent given
the input. The field adds causal memory. The feedback loop makes that
memory \emph{adaptive} (content-dependent, not just positional). This
is the one thing attention provably cannot do and the field provably
can.


% ================================================================
\section{Other Perspectives: Multi-Scale and Hierarchical Fields}
\label{sec:perspectives}

The single-scale, single-layer field presented so far is the minimal
instantiation of stigmergic dynamics.  But the three-ingredient
recipe is a \emph{local} constraint---it specifies the conditions for
edge-of-chaos dynamics within a single field.  This raises a natural
question: what happens when multiple fields interact across timescales
or across layers?  We argue that the uncoupling requirements
(injection and erosion uncorrelated with the field) need only hold
\emph{within} a scale, not \emph{between} scales.  Signals arriving
from a different timescale or a different layer are, from the
perspective of the receiving field, external inputs---structurally
analogous to injection, not to self-correlated feedback.

% ----------------------------------------------------------------
\subsection{Multi-Scale Fields}
\label{sec:multiscale}

Instead of a single field per head, assign each head a
\emph{different} evaporation rate, or equip each head with multiple
fields at distinct timescales.  In the simplest version, head $h$
operates at retention rate $\retain_h$, and the collection of heads
spans a range from fast (small $\retain$, local pattern tracking)
to slow (large $\retain$, global structure).

A richer variant couples a fast field and a slow field per head:
\begin{align}
    \phi_i^{\text{fast}} &= \retain_f \cdot \phi_{i-1}^{\text{fast}} + d_i
    \label{eq:fast_field} \\
    \phi_i^{\text{slow}} &= \retain_s \cdot \phi_{i-1}^{\text{slow}} + h(\phi_i^{\text{fast}})
    \label{eq:slow_field}
\end{align}
where $\retain_s \gg \retain_f$ and $h$ is a learned scalar
projection.  The fast field tracks recent attention patterns via
output deposits as before.  The slow field accumulates a summary of
\emph{which fast-field patterns persist}---it is a distillation of
the distillation.  Both fields contribute additive biases to the
attention logits:
\begin{equation}
    \ell_{ij}^{(h)} = \frac{\q_i^{(h)\top} \K_j^{(h)}}{\sqrt{d_h}}
    + \phi_j^{\text{fast}} + \phi_j^{\text{slow}}
    \label{eq:multiscale_logit}
\end{equation}

The biological analogue is direct: social insects use multiple
pheromone species with different volatilities.  Trail pheromone
(fast) marks recent paths and evaporates in minutes; colony
pheromone (slow) marks established routes and persists for hours.
The slow pheromone does not respond to individual ant decisions---it
integrates over the fast pheromone's history.

\paragraph{Why the three-ingredient recipe still holds.}
Consider the slow field in isolation.  Its deposit source is
$h(\phi^{\text{fast}})$, which depends on the fast field but
\emph{not} on $\phi^{\text{slow}}$ itself.  From the slow field's
perspective, this is injection: an external signal uncorrelated with
its own state.  Its erosion ($\retain_s < 1$) is uniform.  Its
amplification comes indirectly: the slow field biases attention,
which shapes the output, which feeds the fast field, which eventually
feeds the slow field---but this loop passes through the fast field's
own dynamics, maintaining the separation of scales.  The
uncoupling constraints hold \emph{within each scale} even though the
scales are coupled \emph{between} each other.

This is analogous to how weather (slow) affects pheromone evaporation
rates without violating the local dynamics of trail formation: the
weather is external to the ant--pheromone system even though it
modulates it.

% ----------------------------------------------------------------
\subsection{Cross-Layer Fields}
\label{sec:crosslayer}

In a multi-layer transformer, each layer currently maintains an
independent field with no memory of previous layers.  A natural
extension passes field state \emph{vertically}.  For a network with
layers $L = 1, \ldots, \mathcal{L}$, the field at layer $L$ evolves
as:
\begin{align}
    \phi_i^{(L)} &= \retain^{(L)} \cdot \phi_{i-1}^{(L)} + d_i^{(L)},
    \qquad d_i^{(L)} = W_d^{(L)} \bm{o}_i^{(L)}
    \label{eq:crosslayer_recurrence}
\end{align}
with the cross-layer initial condition:
\begin{equation}
    \phi_0^{(L+1)} = g^{(L)}\!\bigl(\phi_T^{(L)}\bigr)
    \label{eq:crosslayer_init}
\end{equation}
where $\phi_T^{(L)}$ is the terminal field state of layer $L$ and
$g^{(L)}$ is a learned scalar projection.  Layer $L+1$ inherits a
landscape shaped by layer $L$'s attention patterns and refines it
with its own deposits.

This creates \textbf{hierarchical feature accumulation}: layer~1's
field summarises token-level patterns (which tokens were attended
to), layer~2's field summarises layer~1's abstractions (which
\emph{patterns of attention} recurred), and so on.  Each layer
refines the previous layer's trail map at a higher level of
abstraction.  Unrolling the full cross-layer field at layer $L+1$:
\begin{equation}
    \phi_i^{(L+1)} = (\retain^{(L+1)})^i \cdot g^{(L)}\!\bigl(\phi_T^{(L)}\bigr)
    + \sum_{j=0}^{i} (\retain^{(L+1)})^{i-j} \cdot d_j^{(L+1)}
    \label{eq:crosslayer_unrolled}
\end{equation}
The first term is the inherited landscape, decaying as the new
layer processes the sequence; the second term is the new layer's own
deposits.  The relative influence of inheritance vs.\ fresh deposits
is controlled by the retention rate and sequence length.

The same within-scale uncoupling argument applies.  Layer $L+1$'s
field receives $g(\phi_T^{(L)})$ as an initial condition---a
one-time injection from an external source.  The field's own
recurrence, deposits, and erosion then proceed independently.  The
cross-layer signal is structurally identical to an initial deposit:
external, uncorrelated with the receiving field's future state.

\paragraph{Combining multi-scale and cross-layer.}
The two extensions compose naturally.  Each layer $L$ maintains fast
and slow fields; the slow field's terminal state seeds the next
layer's slow field:
\begin{align}
    \phi_i^{(L,f)} &= \retain_f^{(L)} \cdot \phi_{i-1}^{(L,f)} + d_i^{(L)}
    \label{eq:combined_fast} \\
    \phi_i^{(L,s)} &= \retain_s^{(L)} \cdot \phi_{i-1}^{(L,s)} + h^{(L)}\!\bigl(\phi_i^{(L,f)}\bigr)
    \label{eq:combined_slow} \\
    \phi_0^{(L+1,s)} &= g^{(L)}\!\bigl(\phi_T^{(L,s)}\bigr)
    \label{eq:combined_cross}
\end{align}
with the attention logit at layer $L$:
\begin{equation}
    \ell_{ij}^{(L,h)} = \frac{\q_i^{(L,h)\top} \K_j^{(L,h)}}{\sqrt{d_h}}
    + \phi_j^{(L,f)} + \phi_j^{(L,s)}
    \label{eq:combined_logit}
\end{equation}
This gives the network a hierarchy along two axes:
\emph{temporal} (fast vs.\ slow within a layer) and
\emph{representational} (shallow vs.\ deep across layers).  The
result is a lattice of interacting fields, each governed locally by
the three-ingredient recipe, coupled only through external channels
that preserve the uncoupling constraints.

% ----------------------------------------------------------------
\subsection{Adaptive Criticality}
\label{sec:adaptive_crit}

A more speculative extension treats the evaporation rate itself as a
field.  A \emph{volatility field} $\psi_i$ modulates the primary
field's retention:
\begin{align}
    \phi_i &= \retain(\psi_i) \cdot \phi_{i-1} + d_i
    \label{eq:adaptive_phi} \\
    \psi_i &= \retain_\psi \cdot \psi_{i-1} + d_i^{(\psi)}
    \label{eq:volatility_field}
\end{align}
where $\retain(\psi) = \sigma(\psi)$ maps the volatility field to a
retention rate via a sigmoid, and $\retain_\psi \approx 1$ ensures
$\psi$ evolves on a much slower timescale than $\phi$.

This directly raises the question of whether the three-ingredient
recipe tolerates inter-scale coupling at all.  The lattice ant model
tested state-dependent erosion and found that it destroys
edge-of-chaos dynamics in every case---but those tests modulated
erosion as a function of the field's \emph{own} value at the
\emph{same} position and timescale.  The multi-scale and cross-layer
extensions proposed above involve a different structure: modulation
by a signal that is external to the receiving field's dynamics.  We
state this distinction as a hypothesis.

\begin{quote}
\textbf{Hypothesis (Scale-local uncoupling).}
\emph{The three-ingredient conditions for edge-of-chaos
dynamics---amplification, uncorrelated injection, and uncorrelated
erosion---are scale-local: they must hold within a single field's
characteristic timescale $\tau_{\emph{eff}}$, but may be violated by
signals arriving from a sufficiently separated timescale.}
\end{quote}

\paragraph{Why this might be correct.}
The argument is an adiabatic approximation.  If the volatility field
$\psi$ varies on a timescale $\tau_\psi \gg \tau_\phi$, then over
any window of length $\sim\!\tau_\phi$ the retention rate
$\retain(\psi)$ is effectively constant.  The fast field experiences
locally uniform erosion at every instant, which is exactly what the
recipe requires.  The slow modulation shifts the fast field's
operating point over time but does not introduce within-scale
correlation between the erosion rate and the field value.  This is
the same reasoning that allows thermodynamics to treat temperature as
a constant in chemical kinetics even though temperature is itself
changing: the reaction timescale is fast relative to the thermal
timescale.  In biological terms, weather modulates pheromone
evaporation rates without violating the local trail-formation
dynamics---rain is external to the ant--pheromone system even though
it changes the evaporation rate everywhere.

The same logic extends to cross-layer fields
(\Cref{sec:crosslayer}): the inherited initial condition
$g(\phi_T^{(L)})$ is a one-time injection from an external source.
And to multi-scale fields (\Cref{sec:multiscale}): the slow field's
deposit $h(\phi^{\text{fast}})$ depends on the fast field, not on
the slow field's own state.  In each case, the inter-scale signal is
structurally external to the receiving field.

\paragraph{Why this might be wrong.}
Three concerns deserve scrutiny:
\begin{enumerate}
    \item \textbf{The ant paper did not test this.}  All seven failed
          mechanisms in \cite{antpaper} involved within-scale
          state dependence.  Slowly varying external modulation was
          not tested.  The hypothesis is an extrapolation, not an
          established result.

    \item \textbf{Bifurcation crossings.}  The adiabatic
          approximation fails when the slow variable passes through a
          bifurcation point of the fast system.  If $\psi$ drifts
          through the critical evaporation rate $\evap^*$, the fast
          field undergoes an abrupt phase transition rather than a
          smooth parameter change.  The fast dynamics cannot track a
          discontinuity adiabatically.  In the neural setting this
          might be benign (the network learns to avoid the critical
          point) or catastrophic (training becomes unstable near the
          transition).  This is an empirical question.

    \item \textbf{``Sufficiently separated'' is vague.}  How large
          must $\tau_\psi / \tau_\phi$ be?  The adiabatic
          approximation is a continuum of accuracy, not a binary
          condition.  At moderate separation the fast field
          experiences erosion that is \emph{approximately} uniform
          but not exactly so---enough correlation to degrade
          edge-of-chaos dynamics without obviously breaking them.
          The required separation ratio is likely task-dependent and
          must be determined empirically.
\end{enumerate}

\paragraph{A testable prediction.}
The lattice ant model provides a direct testbed.  Modulate the
evaporation rate $\lambda$ sinusoidally at period $P$ and measure the
edge-of-chaos order parameter as a function of $P / \tau_\phi$.  The
hypothesis predicts that edge-of-chaos dynamics survive when $P \gg
\tau_\phi$ and collapse when $P \sim \tau_\phi$.  The crossover
point would quantify the required timescale separation.  If
confirmed, this would establish scale-local uncoupling as a general
principle for stigmergic systems; if falsified, it would constrain
the multi-scale and adaptive extensions to purely static
(non-state-dependent) parameter choices.

% ----------------------------------------------------------------
\subsection{Experimental Feasibility at Current Scale}
\label{sec:feasibility}

The micro-GPT used in \Cref{sec:experiments} (2 layers, 4 heads,
$d_{\text{model}} = 48$, $T = 18$) is sufficient to test some of
the proposed extensions but not all.  We summarise the feasibility
constraints to guide the experimental programme.

\paragraph{Multi-scale fields (\Cref{sec:multiscale}).}
\emph{Feasible.}  With 4 heads, the simplest test assigns different
evaporation rates per head---e.g., 2 fast heads ($\evap_f$) and 2
slow heads ($\evap_s$)---and compares against uniform $\evap$ across
all heads.  This requires no architectural changes beyond making
$\evap$ a per-head parameter.  The coupled fast--slow variant
(\Cref{eq:fast_field,eq:slow_field}) adds one scalar projection $h$
per head, a negligible parameter increase.

However, the short sequence length ($T = 18$) constrains the range
of useful timescales.  A slow field with $\retain_s = 0.99$ has
effective memory $\tau_{\text{eff}} = -1/\!\ln(0.99) \approx 100$,
far exceeding the sequence---it is effectively frozen, accumulating
all deposits with near-zero decay.  A fast field with $\retain_f =
0.90$ has $\tau_{\text{eff}} \approx 10$, covering most of the
sequence.  Meaningful separation requires:
\begin{equation}
    1 \ll \tau_{\text{eff}}^{\text{fast}} < T \ll \tau_{\text{eff}}^{\text{slow}}
    \label{eq:timescale_constraint}
\end{equation}
At $T = 18$ there is limited dynamic range for the fast field to
operate in a non-trivial regime while leaving room for a separated
slow field.

\paragraph{Cross-layer fields (\Cref{sec:crosslayer}).}
\emph{Feasible at minimum depth.}  Two layers provide exactly one
cross-layer handoff: layer~1's terminal field state $\phi_T^{(1)}$
seeds layer~2's initial condition via \Cref{eq:crosslayer_init}.
This is the minimal test of whether field inheritance helps at all.
It cannot show whether the effect compounds over depth.  A 4-layer
model provides three cross-layer handoffs---enough to observe whether
hierarchical accumulation compounds or saturates.

\paragraph{Adaptive criticality (\Cref{sec:adaptive_crit}).}
\emph{Not feasible at current scale.}  The volatility field $\psi$
must vary meaningfully \emph{during} the sequence for the adaptive
mechanism to differ from a static $\evap$.  At $T = 18$, a slowly
varying $\psi$ with $\retain_\psi \approx 1$ is effectively constant
across the sequence---indistinguishable from a learned fixed
$\evap$.  Testing adaptive criticality requires sequences long enough
for the slow field to traverse a non-trivial trajectory, suggesting
$T \geq 128$.

\paragraph{Recommended progression.}
The experimental programme proceeds in three stages:
\begin{enumerate}
    \item \textbf{Validate the base mechanism} on a task that requires
          sequential accumulation (state tracking: running sum of
          $\pm 1$ tokens, as proposed in \Cref{sec:right_task}).  If
          the single-scale feedback field does not help here, the
          hierarchical extensions are moot.
    \item \textbf{Test multi-scale and cross-layer} on the same task
          at $T = 18$.  Per-head evaporation rates and one cross-layer
          handoff are cheap additions.  Compare: (a) uniform $\evap$
          feedback, (b) per-head $\evap$ feedback, (c) cross-layer
          feedback, (d) both.
    \item \textbf{Scale up} to $T \in [128, 256]$ with a 4-layer
          model ($d_{\text{model}} \geq 96$) to test adaptive
          criticality and deep cross-layer accumulation.  At this
          scale, the timescale separation
          constraint (\Cref{eq:timescale_constraint}) admits a
          meaningful range, and the bifurcation-crossing concern
          becomes empirically testable.
\end{enumerate}


% ================================================================
\section{Conclusion}
\label{sec:conclusion}

We presented stigmergic attention, an architecture that augments
transformer attention with a decaying pheromone field. Guided by the
phase transition theory of lattice ant models, we identified that the
critical ingredient for non-trivial dynamics is \emph{positive
feedback}---computing deposits from attention output, not input.

The non-feedback variant is provably inert: a causally-weighted
moving average that adds nothing beyond positional information.
The feedback variant closes the loop between field and attention,
creating a nonlinear state-space model coupled to the attention
mechanism.

The length generalization experiments (\Cref{sec:lengthgen}) provide
the clearest evidence for the mechanism. On parity---a canonical
state-tracking task---the feedback field improves out-of-distribution
accuracy from $0.54 \to 0.62$ at $2\times$ the training length and
from $0.51 \to 0.55$ at $4\times$, while all models achieve 1.000
at the training length. The effect is task-specific: it appears on
parity (1-bit state, within the field's capacity) but not on
cumulative sum mod 16 (4-bit state, exceeding scalar capacity).
This pattern suggests a \emph{state-complexity scaling law}: the
field should help when the task's state fits in $O(H)$ scalar
channels.

The results are modest in absolute terms but directionally clear:
the feedback field provides a form of intrinsic memory that
standard attention lacks, and this memory transfers to unseen
sequence lengths because the exponential decay $\retain^{i-j}$
depends on distance, not absolute position. The mechanism maps
cleanly to existing architectural concepts---it is simultaneously a
content-dependent ALiBi, a minimal SSM coupled to attention, and
a neural instantiation of the ant trail recipe.

A fundamental limitation of the scalar-bias architecture is its
expressivity: a single scalar per head can bias \emph{how much}
attention a position receives, but cannot modulate \emph{what kind}
of matching occurs. The field says ``attend more here'' but not
``attend in this direction.'' Real pheromone trails encode
\emph{gradients}---ants follow concentration differentials, not
absolute levels. The neural analogue would be a field that modulates
the geometry of the attention space itself---producing vectors that
reshape keys or queries rather than scalars that nudge logits. This
suggests a successor architecture where the field is a co-equal
routing mechanism, not an additive decoration, decomposing attention
into structured content matching and experiential trail following.
This direction is explored in a companion paper.

\paragraph{Note: from scalar bias to multiscale vector fields.}
Two companion papers develop the successor architecture.  The
dual-channel paper replaces the scalar logit bias with a vector
social field that modulates key geometry, dramatically increasing
state capacity (cumsum mod 16 at $2\times$: 0.999 vs 0.279 baseline).
The hierarchical coupling paper further shows that a single timescale
is insufficient on natural language: adding a second, slow field
with learnable retention produces a 36.6\% ablation delta at late
positions, compared to 2.18\% for the single field, and the model
spontaneously discovers the scale separation that the edge-of-chaos
theory here predicts is necessary---fast field for local dynamics,
slow field for long-range memory, with the evaporation rate
controlling the transition between ordered and disordered phases
exactly as in the lattice ant model.


\bibliographystyle{plain}
\begin{thebibliography}{9}

\bibitem{antpaper}
Pheromone trail formation at the edge of chaos: Phase diagrams of a
lattice ant model.
\emph{Research paper}, 2024.

\bibitem{vaswani2017attention}
A.~Vaswani, N.~Shazeer, N.~Parmar, J.~Uszkoreit, L.~Jones,
A.~N.~Gomez, \L.~Kaiser, and I.~Polosukhin.
Attention is all you need.
In \emph{NeurIPS}, 2017.

\bibitem{gu2022s4}
A.~Gu, K.~Goel, and C.~R\'{e}.
Efficiently modeling long sequences with structured state spaces.
In \emph{ICLR}, 2022.

\bibitem{gu2023mamba}
A.~Gu and T.~Dao.
Mamba: Linear-time sequence modeling with selective state spaces.
\emph{arXiv:2312.00752}, 2023.

\bibitem{sun2023retnet}
Y.~Sun, L.~Dong, S.~Huang, S.~Ma, Y.~Xia, J.~Xue, J.~Wang, and F.~Wei.
Retentive network: A successor to transformer for large language models.
\emph{arXiv:2307.08621}, 2023.

\bibitem{ramsauer2021hopfield}
H.~Ramsauer, B.~Sch\"{a}fl, J.~Lehner, P.~Seidl, M.~Widrich,
T.~Adler, L.~Gruber, M.~Holzleitner, M.~Pavlovi\'{c}, G.~K.~Sandve,
V.~Greiff, D.~Kreil, M.~Kopp, G.~Klambauer, J.~Brandstetter, and
S.~Hochreiter.
Hopfield networks is all you need.
In \emph{ICLR}, 2021.

\bibitem{merrill2023transformers}
W.~Merrill and A.~Sabharwal.
The parallelism tradeoff: Limitations of log-precision transformers.
\emph{Transactions of the Association for Computational Linguistics},
11:531--545, 2023.

\bibitem{press2022alibi}
O.~Press, N.~Smith, and M.~Lewis.
Train short, test long: Attention with linear biases enables input
length extrapolation.
In \emph{ICLR}, 2022.

\bibitem{raffel2020t5}
C.~Raffel, N.~Shazeer, A.~Roberts, K.~Lee, S.~Narang, M.~Matena,
Y.~Zhou, W.~Li, and P.~J.~Liu.
Exploring the limits of transfer learning with a unified text-to-text
transformer.
\emph{JMLR}, 21(140):1--67, 2020.

\bibitem{fan2021feedback}
A.~Fan, T.~Lavril, E.~Grave, A.~Joulin, and S.~Sukhbaatar.
Addressing some limitations of transformers with feedback memory.
\emph{arXiv:2002.09402}, 2021.

\bibitem{lieber2024jamba}
O.~Lieber, B.~Lenz, H.~Bata, G.~Cohen, J.~Osin, I.~Dalmedigos,
E.~Safahi, S.~Meirom, Y.~Belinkov, S.~Shalev-Shwartz, and
O.~Shoham.
Jamba: A hybrid transformer-mamba language model.
\emph{arXiv:2403.19887}, 2024.

\bibitem{de2024griffin}
S.~De, S.~L.~Smith, A.~Fernando, A.~Botev, G.~Cristian-Muraru,
A.~Gu, R.~Harber, L.~Hazan, Y.-H.~Hsu, et~al.
Griffin: Mixing gated linear recurrences with local attention for
efficient language models.
\emph{arXiv:2402.19427}, 2024.

\bibitem{peng2023rwkv}
B.~Peng, E.~Alcaide, Q.~Anthony, A.~Albalak, S.~Arcadinho,
H.~Cao, X.~Cheng, M.~Chung, et~al.
RWKV: Reinventing RNNs for the transformer era.
In \emph{EMNLP Findings}, 2023.

\bibitem{dao2024mamba2}
T.~Dao and A.~Gu.
Transformers are SSMs: Generalized models and efficient algorithms
through structured state space duality.
In \emph{ICML}, 2024.

\end{thebibliography}

\end{document}
